\subsection{Common Source (CS)}\label{sec:amp_cs}

\subsubsection{Overview}
The input drives the gate, the output is taken at the drain and the
source is shared (AC ground). It is the MOS voltage gain stage: inverting,
infinite input impedance at DC and output impedance set by the load.

\subsubsection{Schematic}
\begin{circuitfig}
  \draw (5,0) node[nmos](M){};
  \node[right] at (5.4,0) {$M_1$};
  \draw (0,-2.2) node[ground]{} to[sinusoidal voltage source, l={$V_{in}$}] (0,0) to[european resistor, l={$G_S$}] (2.6,0) node[circ]{} node[below right]{$V_g$} -- (M.G);
  \draw (M.S) -- (5,-2.2) node[ground]{};
  \draw (M.D) -- (5,2) node[circ]{} to[european resistor, l_={$G_D$}] (5,4) node[vcc]{$V_{dd}$};
  \draw (5,2) -- (7.5,2) node[circ]{} node[above]{$V_{out}$} to[capacitor, l={$C_L$}] (7.5,-2.2) node[ground]{};
  \draw (2.6,0) to[capacitor, l={$C_{gd}$}] (2.6,2) -- (5,2);
  \draw (2.6,0) to[capacitor, l_={$C_{gs}$}] (2.6,-2.2) node[ground]{};
\end{circuitfig}
\(G_S\) is the conductance of the driving source; \(C_{gs}\), \(C_{gd}\)
are the device capacitances (Section~\ref{sec:devices_mosfet}) drawn
outside the symbol.

\subsubsection{Transfer Function}
\begin{enumerate}
  \item \textbf{Large Signal, Time Domain, DC Input Analysis (Q5)}: \(I_G = 0\) so \(V_g = V_{in}\).
    KCL at \(V_{out}\):
    \begin{align*}
      (V_{out} - V_{dd})G_D + I_D &= 0 \\
      V_{out} &= V_{dd} - \frac{I_D}{G_D} \tag{1}
    \end{align*}
    \begin{enumerate}
      \item Cutoff, \(V_{in} < V_{th}\): \(I_D = 0\), \(V_{out} = V_{dd}\).
      \item Saturation, \(V_{in} \geq V_{th}\), \(V_{out} \geq V_{in} - V_{th}\):
        \begin{align*}
          V_{out} &= V_{dd} - \frac{k_n (V_{in} - V_{th})^2 (1 + K_{CLM} V_{out})}{G_D} \quad [V_{DS} = V_{out}] \\
          V_{out} &= V_{dd} - \frac{k_n (V_{in} - V_{th})^2}{G_D} - \frac{k_n K_{CLM} (V_{in} - V_{th})^2}{G_D} V_{out} \\
          V_{out}\left(1 + \frac{k_n K_{CLM} (V_{in} - V_{th})^2}{G_D}\right) &= V_{dd} - \frac{k_n (V_{in} - V_{th})^2}{G_D} \\
          V_{out}\left(G_D + k_n K_{CLM} (V_{in} - V_{th})^2\right) &= V_{dd} G_D - k_n (V_{in} - V_{th})^2 \quad [\times G_D] \\
          V_{out} &= \frac{V_{dd} G_D - k_n (V_{in} - V_{th})^2}{G_D + k_n K_{CLM} (V_{in} - V_{th})^2}
        \end{align*}
      \item Edge of saturation (\(K_{CLM} = 0\)), \(V_{out} = V_{in} - V_{th} = V_{ov}\):
        \begin{align*}
          V_{ov} &= V_{dd} - \frac{k_n V_{ov}^2}{G_D} \\
          \frac{k_n}{G_D} V_{ov}^2 + V_{ov} - V_{dd} &= 0 \\
          V_{ov,edge} &= \frac{-1 + \sqrt{1^2 - 4 \cdot \frac{k_n}{G_D} \cdot (-V_{dd})}}{2 \cdot \frac{k_n}{G_D}} \hspace{1cm} [\text{positive root}] \\
          V_{ov,edge} &= \frac{-1 + \sqrt{1 + \frac{4 k_n V_{dd}}{G_D}}}{\frac{2 k_n}{G_D}} = \frac{G_D}{2 k_n}\left(\sqrt{1 + \frac{4 k_n V_{dd}}{G_D}} - 1\right)
        \end{align*}
      \item Triode, \(V_{out} < V_{in} - V_{th}\), with
        \(I_D = K_P \frac{W}{L}\left[(V_{in} - V_{th}) V_{out} - \frac{V_{out}^2}{2}\right] = 2 k_n \left[\ldots\right]\)
        and \(V_x = V_{in} - V_{th}\):
        \begin{align*}
          V_{out} &= V_{dd} - \frac{2 k_n}{G_D}\left[ V_x V_{out} - \frac{V_{out}^2}{2} \right] \\
          V_{out} &= V_{dd} - \frac{2 k_n V_x}{G_D} V_{out} + \frac{k_n}{G_D} V_{out}^2 \\
          \frac{k_n}{G_D} V_{out}^2 - \left(1 + \frac{2 k_n V_x}{G_D}\right) V_{out} + V_{dd} &= 0 \\
          k_n V_{out}^2 - \left(G_D + 2 k_n V_x\right) V_{out} + G_D V_{dd} &= 0 \hspace{1cm} [\times G_D] \\
          V_{out} &= \frac{G_D + 2 k_n V_x - \sqrt{(G_D + 2 k_n V_x)^2 - 4 k_n G_D V_{dd}}}{2 k_n}
        \end{align*}
        taking the smaller root (the larger one gives \(V_{out} > V_x\), outside triode).
    \end{enumerate}
    The amplifier is biased in saturation, between
    \(V_{in} = V_{th}\) and \(V_{in} = V_{th} + V_{ov,edge}\).
  \item \textbf{Small Signal, Frequency Domain, AC Input Analysis (Q4)}: \(V_{dd}\) is AC ground.
    Let \(G_o = G_D + g_o\).\\
    Applying KCL at \(V_g\):
    \begin{align*}
      (V_g - V_{in})G_S + V_g sC_{gs} + (V_g - V_{out})sC_{gd} &= 0 \\
      V_g G_S - V_{in} G_S + V_g sC_{gs} + V_g sC_{gd} - V_{out} sC_{gd} &= 0 \\
      V_g (G_S + s(C_{gs} + C_{gd})) - V_{out} sC_{gd} &= V_{in} G_S \tag{2}
    \end{align*}
    Applying KCL at \(V_{out}\):
    \begin{align*}
      V_{out}(G_D + g_o) + g_m V_g + V_{out} sC_L + (V_{out} - V_g)sC_{gd} &= 0 \\
      V_{out} G_o + V_{out} sC_L + V_{out} sC_{gd} &= V_g sC_{gd} - g_m V_g \\
      V_{out}(G_o + s(C_L + C_{gd})) &= V_g (sC_{gd} - g_m) \\
      V_g &= \frac{G_o + s(C_L + C_{gd})}{sC_{gd} - g_m} V_{out} \tag{3}
    \end{align*}
    Substituting (3) into (2):
    \begin{align*}
      V_{out}\left(\frac{(G_o + s(C_L + C_{gd}))(G_S + s(C_{gs} + C_{gd}))}{sC_{gd} - g_m} - sC_{gd}\right) &= V_{in} G_S \\
      V_{out}\left(\frac{(G_o + s(C_L + C_{gd}))(G_S + s(C_{gs} + C_{gd})) - sC_{gd}(sC_{gd} - g_m)}{sC_{gd} - g_m}\right) &= V_{in} G_S \\
      V_{out}\left[(G_o + s(C_L + C_{gd}))(G_S + s(C_{gs} + C_{gd})) - sC_{gd}(sC_{gd} - g_m)\right] &= V_{in} G_S (sC_{gd} - g_m)
    \end{align*}
    \begin{align*}
      H(s) = \frac{V_{out}}{V_{in}} = \frac{G_S (sC_{gd} - g_m)}{(G_o + s(C_L + C_{gd}))(G_S + s(C_{gs} + C_{gd})) - sC_{gd}(sC_{gd} - g_m)} \tag{4}
    \end{align*}
    Expanding the denominator:
    \begin{align*}
      D(s) &= G_o G_S + s G_o (C_{gs} + C_{gd}) + s G_S (C_L + C_{gd}) \\
           &\quad + s^2 (C_L + C_{gd})(C_{gs} + C_{gd}) - s^2 C_{gd}^2 + s g_m C_{gd} \\
      D(s) &= G_o G_S + s\left[G_o (C_{gs} + C_{gd}) + G_S (C_L + C_{gd}) + g_m C_{gd}\right] \\
           &\quad + s^2 \left[(C_L + C_{gd})(C_{gs} + C_{gd}) - C_{gd}^2\right] \\
      (C_L + C_{gd})(C_{gs} + C_{gd}) - C_{gd}^2 &= C_L C_{gs} + C_L C_{gd} + C_{gd} C_{gs} + C_{gd}^2 - C_{gd}^2 \\
           &= C_L C_{gs} + C_L C_{gd} + C_{gs} C_{gd} \\
      D(s) &= G_o G_S + s\left[G_o (C_{gs} + C_{gd}) + G_S (C_L + C_{gd}) + g_m C_{gd}\right] \\
           &\quad + s^2 \left[C_L C_{gs} + C_L C_{gd} + C_{gs} C_{gd}\right]
    \end{align*}
\end{enumerate}

\subsubsection{Analysis}
\begin{enumerate}
  \item \textbf{DC Gain}: \(s = 0\) in (4),
    \[
      A_{V0} = H(0) = \frac{G_S (0 - g_m)}{(G_o + 0)(G_S + 0) - 0} = \frac{-g_m G_S}{G_o G_S} = -\frac{g_m}{G_D + g_o} = -g_m (R_D \parallel r_o)
    \]
    The largest possible gain (\(G_D \to 0\), current source load) is the
    intrinsic gain \(-\frac{g_m}{g_o} = -\frac{2 I_D / V_{ov}}{K_{CLM} I_D} = -\frac{2}{K_{CLM} V_{ov}}\).
  \item \textbf{Zeros}: \(N(s) = G_S(sC_{gd} - g_m) = 0\):
    \[
      s_z = +\frac{g_m}{C_{gd}} \hspace{1cm} \text{[right half plane: the input feeds forward through } C_{gd}\text{]}
    \]
  \item \textbf{Poles}: Writing \(D(s) = a_0 + a_1 s + a_2 s^2\) and assuming
    widely spaced real poles (Section~\ref{sec:theory_cs_second_order},
    \(p_1 p_2 = \frac{a_0}{a_2}\), \(p_1 + p_2 = \frac{a_1}{a_2}\), \(p_2 \gg p_1\)):
    \begin{align*}
      p_1 &\approx \frac{a_0}{a_1} = \frac{G_o G_S}{G_o (C_{gs} + C_{gd}) + G_S (C_L + C_{gd}) + g_m C_{gd}} \\
      p_1 &\approx \frac{1}{\frac{C_{gs} + C_{gd}}{G_S} + \frac{C_L + C_{gd}}{G_o} + \frac{g_m C_{gd}}{G_o G_S}}
          = \frac{1}{R_S (C_{gs} + C_{gd}) + R_o (C_L + C_{gd}) + g_m R_o R_S C_{gd}} \\
      p_1 &\approx \frac{1}{R_S\left[C_{gs} + C_{gd}\left(1 + g_m R_o\right)\right] + R_o (C_L + C_{gd})} \hspace{1cm} \left[\text{divide by } G_o G_S,\ R_o = \tfrac{1}{G_o}\right] \\
      p_2 &\approx \frac{a_1}{a_2} = \frac{G_o (C_{gs} + C_{gd}) + G_S (C_L + C_{gd}) + g_m C_{gd}}{C_L C_{gs} + C_L C_{gd} + C_{gs} C_{gd}}
    \end{align*}
    The term \(C_{gd}(1 + g_m R_o) = C_{gd}(1 + |A_{V0}|)\) is the Miller
    multiplication of \(C_{gd}\) seen at the gate.
  \item \textbf{Ideal Source} (\(G_S \to \infty\)): divide numerator and denominator of (4) by \(G_S\):
    \begin{align*}
      H(s) &= \frac{sC_{gd} - g_m}{(G_o + s(C_L + C_{gd}))\left(1 + \frac{s(C_{gs} + C_{gd})}{G_S}\right) - \frac{sC_{gd}(sC_{gd} - g_m)}{G_S}}
      \xrightarrow{G_S \to \infty} \frac{sC_{gd} - g_m}{G_o + s(C_L + C_{gd})}
    \end{align*}
    The pole is the root of \(G_o + s(C_L + C_{gd}) = 0\), \(s = -p_1\):
    \begin{align*}
      H(s) &= \frac{sC_{gd} - g_m}{G_o + s(C_L + C_{gd})} \hspace{1cm}
      p_1 = \frac{G_D + g_o}{C_L + C_{gd}} \\
      GB &= \frac{|A_{V0}| p_1}{2\pi} = \frac{1}{2\pi} \cdot \frac{g_m}{G_D + g_o} \cdot \frac{G_D + g_o}{C_L + C_{gd}} = \frac{g_m}{2\pi (C_L + C_{gd})}
    \end{align*}
  \item \textbf{Impedances}:
    \paragraph{At Input:} \(I_G = 0\) at DC, so only capacitances load the gate.
      With \(V_{out} = A_{V} V_g\) at the other end of \(C_{gd}\), the current
      drawn from the gate node is
      \begin{align*}
        I_{in} &= V_g sC_{gs} + (V_g - A_V V_g) sC_{gd} = V_g s\left[C_{gs} + C_{gd}(1 - A_V)\right] \\
        Z_{in} &= \frac{V_g}{I_{in}} = \frac{1}{s\left[C_{gs} + C_{gd}(1 - A_V)\right]} \approx \frac{1}{s\left[C_{gs} + C_{gd}(1 + g_m R_o)\right]} \hspace{1cm} [A_V \approx A_{V0} = -g_m R_o]
      \end{align*}
    \paragraph{At Output:} \(V_{in} = 0\) and (low frequency) \(V_g = 0\), so
      \(g_m V_g = 0\). A test voltage \(V_{test}\) at the output draws
      \begin{align*}
        I_{test} &= V_{test} G_D + V_{test} g_o + g_m \cdot 0 = V_{test}(G_D + g_o) \\
        Z_{out} &= \frac{V_{test}}{I_{test}} = \frac{1}{G_D + g_o} = R_D \parallel r_o
      \end{align*}
  \item \textbf{Swing}: \(V_{ov} \leq V_{out} \leq V_{dd}\) (saturation to cutoff).
\end{enumerate}

\subsubsection{Design}
\begin{enumerate}
  \item Choose \(I_D\) (power) and \(V_{ov}\) (gain vs. swing):
    \(\frac{W}{L} = \frac{2 I_D}{K_P V_{ov}^2}\).
  \item Choose the output bias \(V_{OUT}\) (mid swing):
    \(R_D = \frac{V_{dd} - V_{OUT}}{I_D}\).
  \item Bias the gate at \(V_{IN} = V_{th} + V_{ov}\).
  \item With \(g_m = \frac{2 I_D}{V_{ov}}\), \(G_D = \frac{I_D}{V_{dd} - V_{OUT}}\), \(g_o = K_{CLM} I_D\):
    \begin{align*}
      A_{V0} = -\frac{g_m}{G_D + g_o} &= -\frac{2 I_D}{V_{ov}} \cdot \frac{1}{\frac{I_D}{V_{dd} - V_{OUT}} + K_{CLM} I_D}
              = -\frac{2}{V_{ov}} \cdot \frac{1}{\frac{1}{V_{dd} - V_{OUT}} + K_{CLM}} \\
             &= -\frac{2 (V_{dd} - V_{OUT})}{V_{ov}\left(1 + K_{CLM}(V_{dd} - V_{OUT})\right)} \hspace{1cm} [\times \tfrac{V_{dd} - V_{OUT}}{V_{dd} - V_{OUT}}]
    \end{align*}
\end{enumerate}

\subsubsection{Question}
\(V_{dd} = 1.8\,\)V, \(K_P = 200\,\mu\)A/V\(^2\), \(V_{th} = 0.5\,\)V,
\(K_{CLM} = 0.1\,\)V\(^{-1}\), \(I_D = 100\,\mu\)A, \(V_{ov} = 0.2\,\)V,
\(V_{OUT} = 1\,\)V, \(C_L = 1\,\)pF, ideal source. Find \(\frac{W}{L}\),
\(R_D\), \(A_{V0}\), \(f_{p1}\) and \(GB\) (neglect \(C_{gd}\)).
\paragraph{Solution:}
\begin{align*}
  \frac{W}{L} &= \frac{2 \cdot 100\,\mu A}{200\,\mu A/V^2 \cdot (0.2V)^2} = \frac{200\,\mu A}{8\,\mu A} = 25 \hspace{1cm}
  R_D = \frac{1.8V - 1V}{100\,\mu A} = 8\,k\Omega \\
  g_m &= \frac{2 \cdot 100\,\mu A}{0.2V} = 1\,mS \hspace{1cm}
  g_o = 0.1 \cdot 100\,\mu A = 10\,\mu S \hspace{1cm}
  G_D = \frac{1}{8\,k\Omega} = 125\,\mu S \\
  A_{V0} &= -\frac{1\,mS}{125\,\mu S + 10\,\mu S} = -\frac{1000}{135} = -7.41 \hspace{1cm} (20\log_{10} 7.41 = 17.4\,dB) \\
  A_{V0} &= -\frac{2 \cdot 0.8}{0.2 \cdot (1 + 0.1 \cdot 0.8)} = -\frac{1.6}{0.216} = -7.41 \hspace{1cm} [\text{check with the design formula}] \\
  f_{p1} &= \frac{p_1}{2\pi} = \frac{135\,\mu S}{2\pi \cdot 1\,pF} = 21.5\,MHz \hspace{1cm}
  GB = \frac{1\,mS}{2\pi \cdot 1\,pF} = 159\,MHz
\end{align*}
